\subsection{R-30: tilted padding closes the two-product endpoint range}

\paragraph{Statement R-30.}
Let $n>m\geq3$, let $r\in\{2m-1,2m\}$, and put $a=n-m$.
For every partition $\nu\vdash ns$, $s\geq1$, with exactly
$r$ positive parts and $\nu_1>as$, the two enhanced cones of
G-22 have a common stable point. They have such a point $w$
with the no-scaling property
\[
 g w=cw,\quad g\in\mathrm{SL}_r,\quad c\in\mathbb C^*
       \quad\Longrightarrow\quad c=1.
\]
Thus the common-stability, exact invariant extraction, quotient
 dimension, and all-integer occurrence conclusions of R-29 hold
through $4\leq r\leq2m$ after combining the two results.
The assertion concerns the enhanced tensor: at these endpoints
we do not assert that the underlying padded form itself has
trivial projective stabilizer.

\paragraph{Background.}
Use the tensor-induced metrics and full moment matrix $M$
of R-29. Besides the cited Kempf--Ness criterion, use its
stabilizer consequence at moment zero: the compact stabilizer
Lie algebra is a real form of the complex stabilizer Lie algebra.
The exact inspected locator is Losev (2006), proof of Theorem~1,
printed p.~2, equations~(3)--(4) and the sentence immediately
following (G-B14). In particular a moment-zero vector with
finite compact stabilizer has finite complex stabilizer.
We use this cited fact only for the full ordered padded tensor,
not for its unbalanced polynomial factors.

\begin{proof}
We extend the product normalization at a nonconstant endpoint,
then tilt the padding line so that neither endpoint block has
a constant spectrum.

\paragraph{1. Nonconstant spectra in exactly $m$ variables.}
Let $b_1,\ldots,b_m>0$ be rational, with sum $m$, and suppose
they are not all equal. There is a unit product $P$ of $m$
independent linear forms with
\[
 M(P)=\operatorname{diag}(b_1,\ldots,b_m)
\]
whose projective stabilizer in the unitary group is finite.
Its complex projective stabilizer still contains a torus;
no contrary assertion is needed.

To prove this, sort the $b_i$, clear denominators with $d$,
and put $\lambda=db$ and $\theta=\lambda-\lambda_m\mathbf1$
as in R-29. Now $\theta\ne0$. Take the same Fourier vectors
$u_j$ for $0\leq j<m$; they form a basis. The ordered tensor
\[
 \xi=(u_0\otimes\cdots\otimes u_{m-1})^{\otimes d}
                                      \otimes v_\theta
\]
has moment zero. Its projective stabilizer is finite, as follows.
Write the basis matrix as $DF$, with $D$ positive diagonal and
$F$ the Fourier matrix. An operator preserving each ordered
factor line has the form
\[
 g=DF\operatorname{diag}(t_0,\ldots,t_{m-1})F^{-1}D^{-1}.
\]
Since $\theta\ne0$ and $\theta_m=0$, preserving the auxiliary
extremal line preserves at least one nonzero proper standard
dual flag subspace. On the original space this is a proper
consecutive coordinate subspace. The diagonal $D$ preserves
that subspace, so the circulant matrix
$C=F\operatorname{diag}(t)F^{-1}$ preserves it as well.

Write $C=\sum_{h=0}^{m-1}c_h S^h$ for the cyclic shift $S$.
A nonempty proper consecutive interval in $\mathbb Z/m\mathbb Z$
is not invariant under any nonzero translation: its unique
transition from outside to inside would otherwise have a
nonzero translation period. For each $h\ne0$ there is therefore
an interval index mapped outside by $S^h$. The corresponding
matrix entry of $C$ is $c_h$, and preservation forces $c_h=0$.
Thus $C$ and $g$ are scalar. In $\mathrm{SL}_m$ only finitely
many scalars remain. This proves the asserted finite projective
stabilizer of $\xi$.

Kempf--Ness and the every-subgroup minimum-weight argument in
R-29 show that $P_0^{\otimes d}\otimes v_\theta$, where
$P_0=\prod_j u_j$, is stable. The same reverse Kempf--Ness and
unitary partial-flag normalization give a product $P$ with the
prescribed matrix. To check the new stabilizer assertion, let
$U$ be a unitary projective symmetry of $P$. Moment equivariance
implies that $U$ commutes with $\operatorname{diag}(b)$.
It therefore preserves its eigenspace decomposition and the
extremal line of $v_\theta$. Hence $U$ is also a projective
symmetry of the stable tensor $P^{\otimes d}\otimes v_\theta$.
Modulo scalar unitary matrices, this group is finite.
This proves the endpoint lemma. For block dimensions at most
$m-1$, the stronger complex-stabilizer statement of R-29
remains available.

\paragraph{2. A rational two-eigenvalue tilt.}
Put $\delta_i=\nu_{i+1}/s$ for $0\leq i<r$, so that
$\delta_0>a$ and $\sum_i\delta_i=n$. There is an index $j>0$
with $\delta_j<\delta_0$. Otherwise $\delta_0=n/r$, and
$r>m$ and $a\geq1$ would give
\[
 \delta_0=\frac{m+a}{r}\leq\frac{m+a}{m+1}\leq a,
\]
a contradiction.

Set $A=\delta_0-a>0$ and $B=\delta_j>0$. Choose a sufficiently
small positive rational $\epsilon$ and put
\[
 \alpha=\max(A,B)+\epsilon,\qquad
 \gamma=\min(A,B)-\epsilon.
\]
Require $\gamma>0$ and require both $\alpha,\gamma$ to differ
from every $\delta_i$ with $i\notin\{0,j\}$. Only finitely
many choices of $\epsilon$ are excluded. Since
$\alpha\gamma<AB$, the rational number
\[
 t=\frac{AB-\alpha\gamma}{a(\delta_0-\delta_j)}
\]
lies in $(0,1)$ when $\epsilon$ is small enough. Take the unit
padding line
\[
 \ell=\sqrt{1-t}\,x_0+\sqrt t\,x_j.
\]
The positive Hermitian matrix
\begin{equation}\label{R30:C}
 C=\operatorname{diag}(\delta_0,\ldots,\delta_{r-1})
                                      -a\ell\ell^*
\end{equation}
has eigenvalues $\alpha,\gamma$ in the $(0,j)$ plane and
$\delta_i$ elsewhere. Indeed the plane trace is $A+B$
and its determinant is
\[
 AB-a(\delta_0-\delta_j)t=\alpha\gamma.
\]
All its eigenvalues are positive rational numbers and
$\operatorname{tr}C=m$. The two newly chosen eigenvalues
are distinct from one another and from all the others.

\paragraph{3. Finite compact stabilizers suffice at moment zero.}
Split an orthonormal eigenbasis of $C$ into blocks of dimensions
$k_1=\lfloor r/2\rfloor$ and $k_2=\lceil r/2\rceil$,
putting $\alpha$ in the first block and $\gamma$ in the second.
Both block sizes lie between two and $m$, and each block
contains at least one other eigenvalue. Its spectrum is
therefore nonconstant. Let $w_b$ be its eigenvalue sum divided
by $m$, and divide its eigenvalues by $w_b$. They are positive
rational numbers with sum $m$.

Use paragraph~1 when $k_b=m$, and R-29 when $k_b<m$, to
choose a unit product in each block with this normalized
moment. Call them $P,Q$. Both have finite unitary projective
stabilizers and independent constant directional derivatives.
For
\[
 q=\sqrt{w_1}P+\sqrt{w_2}Q
\]
the orthogonal-block calculation in R-29 gives $M(q)=C$.
The same mixed-bidegree calculation, now restricted to
skew-Hermitian infinitesimal matrices, proves that the unitary
projective stabilizer of $q$ is finite: the off-diagonal blocks
vanish because $m\geq3$ and the derivatives are independent;
the two diagonal blocks are scalar by the compact stabilizer
statements, and their scalars agree because both forms have
degree $m$.

Let $\eta=\nu-\nu_r\mathbf1$ and take the standard unit
extremal vector $v_\eta$. The ordered tensor
\[
 \widetilde\xi=(\ell^{\otimes a}\otimes q)^{\otimes s}
                                             \otimes v_\eta
\]
has full moment
\[
 sa\ell\ell^*+sC-\operatorname{diag}(\eta)
                                               =\nu_r I_r.
\]
Its compact stabilizer is finite because a projective symmetry
of this pure tensor preserves the line of its $q$ factor.
The cited real-form consequence of Kempf--Ness then makes
its complex vector stabilizer finite. Its complex projective
stabilizer has the same Lie algebra: if $B\widetilde\xi
=c\widetilde\xi$ for $B\in\mathfrak{sl}_r$, the zero moment
implies $\langle B\widetilde\xi,\widetilde\xi\rangle=0$
by complex linearity, hence $c=0$. Thus $\widetilde\xi$ is
stable. Multiplication preserves every one-parameter-subgroup
minimum weight, so $(\ell^a q)^{\otimes s}\otimes v_\eta$
is stable. The two-product cycle and block-padding constructions
place it in both enhanced cones.

\paragraph{4. A generic auxiliary flag removes the remaining finite group.}
We now use the irreducible parameter family
\[
 q=\prod_{i=1}^mL_i+\prod_{i=1}^mM_i,\qquad F=\ell^a q,
\]
with all factors and $\ell$ free. The finite-projective-stabilizer
locus for $F$ is a nonempty open in this family. We give explicit
witnesses to establish nonemptiness.

For $r=2m$, take two independent coordinate blocks $x_1,\ldots,x_m$
and $y_1,\ldots,y_m$, let
\[
 P=\prod_{i=1}^m x_i,\qquad Q=\prod_{i=1}^m y_i,
                    \qquad \ell=\sum_i x_i+\sum_i y_i.
\]
For $r=2m-1$, take $m-1$ independent $x$ variables and $m$
independent $y$ variables, and instead set
\[
 P=x_1\cdots x_{m-1}(x_1+\cdots+x_{m-1}),\qquad
 Q=\prod_{i=1}^m y_i,\qquad
 \ell=\sum_i x_i+\sum_i y_i.
\]
In the latter case the $m$ factors of $P$ form a projective frame,
so its projective stabilizer is finite. In both cases $q=P+Q$
is irreducible: as a polynomial in $y_1$ it is primitive and
linear, with coprime leading and constant coefficients
$y_2\cdots y_m$ and $P$. It is irreducible over the fraction
field of the remaining variables, hence by Gauss over their
polynomial ring.

Any infinitesimal projective symmetry of $F$ preserves the
unique nonlinear irreducible factor $q$ and the padding line.
The mixed-bidegree argument makes a symmetry of $q$ block
 diagonal. At $r=2m$ its two blocks are diagonal in the factor
bases; preserving the displayed $\ell$ forces all diagonal
entries to be the same. At $r=2m-1$ the first block is scalar,
and the second is diagonal; the same padding-line condition
again forces one common scalar. Thus the projective stabilizer
Lie algebra of $F$ is zero, proving the required finiteness.

The enhanced stable locus is nonempty by paragraph~3 and open.
Intersect it with the finite-projective-stabilizer open for $F$,
allowing the auxiliary extremal line to vary in its partial
flag orbit. This is one irreducible parameter space, so the
intersection is nonempty. Fix an $F$ in it and let $H$ be its
finite projective automorphism group. Here $\eta\ne0$: if
$\nu$ were rectangular, the same inequality as in paragraph~2
would contradict $\nu_1>as$.

The action of $\mathrm{PGL}_r$ on the nontrivial partial flag
orbit is faithful. Indeed any operator fixing every subspace
of one proper nonzero dimension fixes every line, by intersecting
such subspaces, and is scalar. Consequently each nonidentity
$h\in H$ has a proper closed fixed locus on this orbit. The
stable flags form a nonempty open in the fixed-$F$ fiber.
Choose a flag there outside the finite union of these fixed
loci. If $g w=cw$ at the resulting enhanced point, $g$ must
preserve both $[F]$ and this flag; its projective class is
therefore the identity. Thus $g$ is central in $\mathrm{SL}_r$.
The central exponent is $ns-|\eta|=r\nu_r$, so it acts
trivially and $c=1$.

The exact extraction and dimension arguments of G-22, followed
by the closed-orbit separation and numerical-semigroup argument
of R-28, now give the stated consequences and every integer
stretch $j\geq\beta^2$ for a common ambient generator bound.
\end{proof}

\paragraph{Scope.}
This proof uses finite compact stabilizers at a balanced full
tensor, not a false finite-complex-stabilizer assertion for a
product in $m$ variables. At the final step it removes the
finite padded-form group with the nontrivial auxiliary flag;
it does not claim that a generic two-product padded form is
projectively free. The present proof stops at $2m$ rows. The two sets of $m$
factors span at most $2m$ variables, although an independent
padding line can add one more essential variable. Whether
that additional variable admits the same strict-direction
stability is a separate question; no obstruction beyond the
proved range is asserted.
